% Campaign digest, 2026-09-11/12. Standalone: pdflatex campaign_digest_20260912.tex
% Separate from the three result tables (eval_table_*, latency_table_*, interleaved_table_*);
% this collects what the two days ESTABLISHED, including the negative results.
\documentclass[10pt]{article}
\usepackage[margin=0.9in]{geometry}
\usepackage{booktabs,amsmath,xcolor}
\usepackage[colorlinks=true,linkcolor=black,urlcolor=blue]{hyperref}
\newcommand{\hh}[1]{\textbf{#1}}
\title{AppCorr campaign digest --- 2026-09-11/12}
\author{Qwen3.5 (4B / 35B-A3B / 122B-A10B-FP8), served vLLM 0.28 + HF in-process}
\date{Generated 2026-09-12}
\begin{document}\maketitle

\section*{1. What was built: the unified (vision $+$ decoder) schedule}

One depth axis of $27$ vision-tower stages followed by $L$ decoder layers, cut into $g$ rounds of
equal cumulative cost, so the tower's approximate pass is itself spread over the rounds and the LLM
opens only when the frontier crosses the projector. On 35B V*Bench ($N{=}3341$) the $g{=}4$ bounds
are $[11,22,41,67]$: tower depths $11/22/27/27$ and decoder depths $0/0/13/40$, i.e.\ only two of
the four rounds reach the LLM. $k{<}1$ selects each band by the layer-mean attention of the stages
walked \emph{so far} (progressive), so its $k{<}1$ cells are not a schedule-only comparison.

Three engine changes were needed to make the served form match the schedule. Each is gated and
each is measured; the ladder below is 35B V*Bench Crit.\ Lat.\ (ms, $d{=}150$ band spacing).

\begin{table}[h]\centering\small
\begin{tabular}{lccc l}
\toprule
engine form & $k{=}1$ & $k{=}0.5$ & $k{=}0.25$ & what changed \\
\midrule
unified, as first implemented        & 157 & 110 & 95   & --- \\
$+$ \texttt{correct\_rows} after the crossing & 153 & 123 & 105 & bitwise-identical messages; no net win \\
$+$ \texttt{open\_walk}               & 80  & 56  & 36.5 & engine skips its stock full-depth prefill \\
$+$ \texttt{CONTIG\_ROWS}             & \hh{61.6} & \hh{48.6} & \hh{36.3} & contiguous rows as ONE prefill sequence \\
\midrule
depth-staged, same engine            & 62.0 & 48.5 & 36.2 & reference \\
interleaved (unstaged), same engine  & 61.5 & 49.8 & 36.4 & reference \\
\bottomrule
\end{tabular}
\end{table}

\noindent\textbf{Decomposition that drove it.} Crit.\ Lat.\ $=$ (wait for the engine) $+$ (last
band's vision correct) $+$ (LLM correct step). The first form spent $65$\,ms waiting because the
engine ran a full-depth prefill of 3341 tokens at the projector crossing and was still busy when
the last band arrived. \texttt{correct\_rows} shortened the vision correct but the wait grew by the
same amount --- the two are a \emph{max}, not a sum, which the first decomposition got wrong.
\texttt{open\_walk} removed the wait; \texttt{CONTIG\_ROWS} cut the engine's 3294-row frontier walk
from 132\,ms to 42.5\,ms by submitting contiguous rows as one prefill sequence instead of $|P|$
one-token pseudo-sequences.

\section*{2. What the unified schedule buys}

\begin{table}[h]\centering\small
\begin{tabular}{llcccc}
\toprule
 & $k$ & streaming & interleaved & depth-staged & \hh{unified} \\
\midrule
\multicolumn{6}{l}{\emph{Total compute, \% of a full-resolution pass (35B)}}\\
V*Bench     & 1.00 & 162 & 195 & 183 & \hh{161} \\
            & 0.50 & 130 & 147 & 141 & \hh{130} \\
            & 0.25 & 115 & 123 & 120 & \hh{115} \\
RealWorldQA & 1.00 & 152 & 197 & 181 & \hh{162} \\
            & 0.50 & 126 & 149 & 141 & \hh{132} \\
\midrule
\multicolumn{6}{l}{\emph{122B V*Bench, Crit.\ Lat.\ (ms) and total compute}}\\
Crit.\ Lat. & 1.00 & 86 & 85 & 202 & 89 \\
Comp.       & 1.00 & 133\% & 201\% & 177\% & \hh{165\%} \\
\bottomrule
\end{tabular}
\end{table}

\noindent\textbf{Baseline provenance (added 2026-09-12, after a correction).} The streaming
Crit.\ Lat.\ figures above and in Table~\ref{tab:interleaved_results} were originally measured
2026-09-09, before three engine changes, while the interleaved and unified figures were measured
after them. Re-probing every streaming cell on the current engine (same server flags, same
$36$ samples, two passes) moved 35B Crit.\ Lat.\ by $-9.1\%$ on average, $32$ of $33$ cells
negative; 122B moved $-2.5\%$, $14$ of $18$ negative. The consequence differs by model. Comparing
interleaved against streaming cell by cell on the $30$ / $24$ cells measurable in both eras:

\begin{center}\small
\begin{tabular}{lccc}
\toprule
 & interleaved faster & streaming faster & tied ($<0.5$\,ms) \\
\midrule
35B, stale baseline & 29 & 1 & 0 \\
35B, one session    & 12 & 14 & 4 \\
122B, stale baseline & 19 & 3 & 2 \\
122B, one session    & 17 & 4 & 3 \\
\bottomrule
\end{tabular}
\end{center}

\noindent So the interleaved schedule's Crit.\ Lat.\ advantage on 35B was an artefact of a stale
streaming baseline and does not survive one-session measurement; on 122B it survives essentially
unchanged. The mechanism is the one the table caption already states: the correct step is GPU-bound
with a floor set by streaming the MoE expert weights once per layer, not by token count, so at the
35B's $19$--$60$\,ms absolute times the saved tokens are swamped, while at the 122B's
$32$--$307$\,ms they are not. Crit.\ Comp.\ is a closed form over the same rows and is unaffected
by any of this.

\noindent\textbf{Reading.} Unified reaches \emph{streaming's} total compute while keeping
\emph{interleaved's} critical compute ($k/g$ of the last round), at a Crit.\ Lat.\ equal to the
other interleaved forms. It cannot beat them on Crit.\ Lat.: its critical compute is identical to
depth-staged by construction. End-to-end it finishes 100--210\,ms earlier on 35B V*Bench because
the tower's approximate work overlaps the arrival gaps. Accuracy is a tie against depth-staged on
all six paired cells (every bootstrap CI covers 0); two independent served runs agreed per sample.
The one place depth-staged loses badly is 122B V*Bench, where each of its four rounds exceeds the
150\,ms band gap on a 3.3\,k-token prompt and the last band queues behind round three.

\section*{3. Measurement-validity findings (the expensive half of the campaign)}

\begin{table}[h]\centering\small
\begin{tabular}{p{0.3\textwidth}p{0.62\textwidth}}
\toprule
finding & evidence \\
\midrule
\hh{The degrade filter was mismatched within table rows.} Bounds measured 2026-08-31 used
\emph{bicubic}; every arm from 09-09 on used \emph{box}. &
Re-running ChartQA's floor under bicubic on the new box gives \hh{60.80} vs the archived
\hh{60.76} (105 discordant, 53/52, $p{=}1.00$); box gives \hh{46.64}. So the $-14.12$\,pp floor
move is the filter, not a regression. Controls: ceiling never consumes the degraded base and moved
$+0.16$\,pp (1.1\% discordance); TextVQA, \emph{pyr on both sides}, moved $+0.24$\,pp (2.0\%) ---
an 11-fold difference in floor discordance between the dataset where the filter was held fixed and
the one where it was not. \\
\addlinespace
\hh{Convention fixed (user decision).} box for realworldqa / chartqa / mmvp / vsr / cvbench;
pyr for textvqa / vstar / visdrone\_* / refcoco / infovqa. &
Per-dataset-family \emph{by decision}. All 18 35B bounds arms re-measured on one box with current
code; superseded files archived with a provenance note. Never pair arms measured under different
filters. \\
\addlinespace
\hh{On-box runs are bit-deterministic; cross-box runs are not.} &
A fresh run of ChartQA k0.50 reproduced the landed arm \hh{byte-identically} (0/2500 discordant);
confirmed twice more (RealWorldQA 765, VSR 1222 --- the latter 0 discordant inside the rebounds run
itself). Across box/env/code with the filter held fixed: $\sim$1--2\% of rows flip with no net
direction on the simple arms, up to $-1.32$\,pp ($p{=}0.00071$) on the streaming arm, which carries
the correction machinery. \\
\addlinespace
\hh{Four significance verdicts changed once rows shared a box, filter and code.} &
MMVP k0.50 vs floor $p$ $0.332 \to 0.031$; MMVP k0.25 no longer below floor (that anomaly was the
artefact); RealWorldQA k0.50/k0.25 $p$ $0.098/0.070 \to 0.033/0.005$; CV-Bench k0.50 $0.037 \to
0.013$. \hh{VSR stayed null in both directions} --- which is what shows the procedure removed an
error rather than manufactured significance. \\
\addlinespace
\hh{The deferred patch score costs a little accuracy.} It is a throughput knob, not a bug. &
Same box/env/code, only \texttt{pscore\_defer} varying: ChartQA $-0.28$\,pp ($p{=}0.392$),
RealWorldQA $-1.44$\,pp ($p{=}0.0127$); pooled 24 better / 42 worse over 66 discordant,
$p{=}0.0356$. The gate shows the implementation is exact: bands 1--3 select identical groups
(Jaccard $1.0000$), the attention term is bitwise equal, and $t_{\text{open}}$ drops 21--56\%. \\
\addlinespace
\hh{Arms can skip different rows; scoring each on its own survivors biases the row.} &
122B InfoVQA: all arms share 338 \texttt{prompt\_too\_long} skips, but the $k{<}1$ interleaved arms
add 194--236 \texttt{oom} skips. Restricting every arm to the common 2227 rows moves floor
$-1.44$\,pp while the $k{<}1$ arms move $-0.3$ to $-0.4$. Every arm is now scored on the
intersection of all arms' scored rows, so the table is paired by construction. \\
\bottomrule
\end{tabular}
\end{table}

\subsection*{3.1 Saturation is a (model, dataset) property}

A row whose floor and ceiling are statistically indistinguishable carries no information about the
method. Model size does \emph{not} predict it.

\begin{table}[h]\centering\small
\begin{tabular}{lccl}
\toprule
dataset & 35B span ($p$) & 122B span ($p$) & \\
\midrule
VSR            & $+0.82$ (0.268)      & $-0.65$ (0.403)      & \hh{saturated on both} \\
RealWorldQA    & $+3.40$ (0.0073)     & $+1.05$ (0.434)      & \hh{saturated on 122B only} \\
MMVP           & $+5.00$ (0.0026)     & $+10.00$ (1.4e-06)   & 122B gap is DOUBLE \\
CV-Bench       & $+1.10$ (0.0172)     & $+1.44$ (0.00060)    & 122B slightly wider \\
ChartQA        & $+42.08$ (7e-270)    & $+36.48$ (6e-217)    & 122B \emph{narrower} \\
V*Bench        & ---                  & $+10.47$ (3.6e-05)   & real headroom \\
VisDrone Count & ---                  & $+5.62$ (1.4e-17)    & real headroom \\
\bottomrule
\end{tabular}
\end{table}

\noindent From 35B to 122B the gap widens on two datasets, narrows on one, and stays saturated on
one --- non-monotone in both directions. Both shortcuts fail: inheriting saturation from the
dataset, and predicting it from model scale. \hh{Test the cell.} VSR is excluded from the
interleaved table on exactly this ground; the remedy for a saturated cell is a harsher degradation
level (L3).

\section*{4. Results across model scale}

\begin{table}[h]\centering\small
\begin{tabular}{llccccc}
\toprule
 & dataset & floor & Ours 25\% & Ours 50\% & $k{=}1$ & ceiling \\
\midrule
\multicolumn{7}{l}{\emph{Recovery of the floor-to-ceiling gap --- the quantity that transfers across scales}}\\
35B & TextVQA (17.31\,pp) & --- & 77.0\% & 89.4\% & 95.3\% & --- \\
4B  & TextVQA (19.34\,pp) & --- & 71.3\% & 85.8\% & 94.3\% & --- \\
\midrule
\multicolumn{7}{l}{\emph{122B, full split, $k{=}0.5$ recovery}}\\
    & ChartQA  & 52.04 & 81.52 & 86.72 & 88.52 & 88.52 \quad(95\%) \\
    & CV-Bench & 84.80 & 85.63 & 86.24 & 86.47 & 86.24 \quad(100\%) \\
    & MMVP     & 78.00 & 79.67 & 80.00 & 85.33 & 88.00 \quad(20\%) \\
    & VSR      & 90.75 & 91.33 & 91.90 & 90.75 & 90.10 \quad(saturated) \\
\bottomrule
\end{tabular}
\end{table}

\noindent\textbf{Two readings that matter.} (i) An $8\times$ model-size difference moves the
\emph{recovered fraction} by 1 point at $k{=}1$ and 4--6 points at reduced keeps, while the raw
accuracies differ by 3--5\,pp and only restate that the bigger model is better. Report recovery.
(ii) On 122B MMVP the keep arms are significantly \emph{below} ceiling ($p{=}7$e-05) and not
distinguishable from floor --- but at $n{=}300$ that is a failure to \emph{demonstrate} benefit,
not a demonstration of its absence: CV-Bench's smaller $+1.44$\,pp reached $p{=}4.1$e-05 on 2638
rows. Every well-powered 122B row recovers 95--100\% at $k{=}0.5$; the one that looks bad is the
one with the least power. \hh{Always print the recovery \% beside the sample size.}

\section*{5. Hypotheses proposed and falsified (five, all within hours, by design)}

\begin{table}[h]\centering\small
\begin{tabular}{p{0.27\textwidth}p{0.30\textwidth}p{0.33\textwidth}}
\toprule
hypothesis & discriminating test & outcome \\
\midrule
The deferred score hurts because band 0 is chosen on residual energy alone, without the attention
term. &
Gate on an affected dataset AND a clean control (VSR). &
\hh{Falsified.} VSR's band-0 selection diverges MORE (Jaccard 0.736/0.594 vs ChartQA's
0.848/0.668) and it pays no accuracy penalty at all. Magnitude of selection change does not
predict, and here inversely tracks, the cost. \\
\addlinespace
The unstaged interleaved schedule degrades progressively with keep, because unselected rows are
never revisited while staged/unified carry every row through the next layer band. &
The deficit must be \emph{monotone}: absent at $k{=}1$, present at $k{=}0.5$, larger at $k{=}0.25$.
Rows already on disk. &
\hh{Falsified,} then the effect itself \hh{closed by pre-registration.} Monotonicity: $+0.47$
(null) / $-1.36$ ($p{=}0.00021$) / $-0.51$ (null) --- the same discordant count (72) either side of
the significant point, so not a power artefact. The cell was then pre-registered (single cell,
direction fixed, decision rule fixed, no re-scan) and tested on five further datasets:
RealWorldQA 30 disc.\ $p{=}0.856$ (adequate); VisDrone Det 19 disc.\ $p{=}0.648$ (borderline);
InfoVQA 155 disc.\ $p{=}0.148$ (adequate$+$); ChartQA 73 disc.\ $p{=}0.101$ \emph{with the sign
reversed} (adequate); and the decider, \hh{122B TextVQA, 163 discordant pairs --- more than double
the originating dataset's 72 --- $-0.26$\,pp, $p{=}0.347$.} Two further attempts (V*Bench 0
discordant, 35B MMVP 4) had \emph{no power} and are excluded rather than counted.
\hh{Five non-replications, zero replications. Reported as no effect}, with the original
significant point left unexplained rather than explained away. \\
\addlinespace
The 122B OOMs come from $|P|$ one-token pseudo-sequences exceeding \texttt{max\_num\_seqs}, so
lowering it to 256 is the fix. &
$|P|$ scales with keep, so $k{=}0.25$ must OOM on strictly fewer rows than $k{=}0.5$. &
\hh{Falsified.} The OOM sets are \emph{bit-identical} across keeps (194/194 unstaged, 236/236
staged on InfoVQA; $\{62,153\}$ both keeps on 122B V*Bench, independently). The failing allocation
is keep-independent --- the per-token side buffer, plus the frontier walk for staged, which
explains why staged's set is a strict superset. \hh{\texttt{max\_num\_seqs} is the wrong knob;}
\texttt{--gpu-mem} is the only one that helps. \\
\addlinespace
Bigger models saturate more rows, so the benchmark suite stops being able to measure the method. &
Enumerate every (model, dataset) span, not only the ones that prompted the sentence. &
\hh{Falsified} by the table in \S3.1 --- non-monotone in both directions, with MMVP (the 122B gap
is double the 35B's) as the direct counterexample. \\
\bottomrule
\end{tabular}
\end{table}

\noindent\textbf{A fifth, caught by the enumeration rule itself.} ``Grounding tasks separate the
$k{=}0.5$ and $k{=}0.25$ arms; VQA tasks mostly do not'' --- prompted by RefCOCO ($p{=}0.0064$) and
VisDrone Count ($p{=}0.0042$). Enumerating all 16 rows across three models kills it: VQA separates
on 6 of 13 and holds the four strongest separations in the campaign (4B TextVQA $p{=}3.1$e-14, 122B
ChartQA $1.1$e-20, 35B ChartQA $5.2$e-11, 35B TextVQA $1.8$e-09). Task type was \emph{confounded
with $n$}: the grounding rows that prompted the sentence are RefCOCO (8811) and VisDrone Count
(2350), two of the largest in either table, while the one small grounding row (VisDrone Det, 448)
fails to separate like the small VQA rows. What actually predicts separation is paired evidence:
every row above ${\sim}100$ discordant pairs separates, every row below ${\sim}30$ does not,
regardless of task or model. The defensible sentence is therefore about the measurement, not the
method: \hh{the two keep levels are distinguishable wherever there is enough paired evidence to
distinguish them.}

\medskip
\noindent\textbf{The transferable guard.} In every case the disconfirming measurement was already
on disk when the general sentence was written. A mechanism makes a result \emph{feel} explained,
and that is exactly when it stops being checked. Write the discriminating prediction down first,
then run it; two points cannot establish a trend, three can test one; and before any sentence of
the form ``$X$ tends to $Y$'', enumerate the cases it covers and compute all of them, including the
ones that did not prompt it.

\section*{6. Campaign state}

Complete: the 35B main-table block (12 keep arms $+$ the missing VSR $k{=}1$, all bounds
re-measured on one box/filter/code); the interleaved table's original six datasets on both models;
the 122B interleaved block. In flight: 122B full-split replacement of the strided subsets
(10.5\% $\to$ 100\% of rows); the interleaved table's extension to ChartQA, CV-Bench, MMVP and
RefCOCO on both models; the 4B main-table block. Parenthesised cells are reduced-$n$ subsets,
flagged per row on \emph{scored} coverage rather than file size.

\end{document}
